\documentclass{article}

\usepackage{float}
\restylefloat{table}

\usepackage{booktabs}
\usepackage{hyperref}

\title{Team Productivity: POC\\\progname}

\author{\authname}

\date{}

\input{../Comments.text}
\input{../Common.text}

\begin{document}

\maketitle

\noindent This document summarizes the contributions of each team member up to 
the POC
Demo.  The time period of interest is the time between the beginning of the term
and the POC demo.\\

\noindent Project URL: 
\href{https://github.com/Aeolus-Software/RealmForge}{RealmForge 
Repository}

\section{Demo Plans}

PoC Demonstration Plan (Taken from the project's Development Plan):\\

\noindent The main risks to the project's development that were identified 
include a 
\textbf{Large Technical Learning Curve}, the possibility of \textbf{Scope 
	Creep}, and a \textbf{Failure to Implement Variation Capabilities}.
\\

\noindent \textbf{Large Learning Curve:} Each project team member has varying 
experience in the field of game development, and some team members may have 
little experience with the game engine used to develop the project's software. 
Team members who are not experienced with the game engine used in the project 
may suffer from a large learning curve, delaying their ability to contribute to 
the project software's source code while they improve their skills. This could 
lead to a failure of the team to complete the development of the software 
within course deadlines.
\\ 

\noindent \textbf{Scope Creep:} Each project team member is passionate about 
video games and game development, and has many ideas for how the project's 
software (the RealmForge game) could function, behave, and look. However, the 
software must be developed within the bounds of the time available in the 
Capstone course. Scope creep could occur if the planned scope for the software 
became larger than what is feasible within the Capstone course's constraints, 
if the team plans to implement a large amount of features from their various 
ideas. This could lead to a failure of the team to complete the development of 
the software within course deadlines, as the majority of development time 
would be spent on individual complex components of the entire software.
\\

\noindent \textbf{Failure to Implement Variation Capabilities:} The purpose of 
the project is to develop a video game, but specifically to develop a video 
game with explicit feature variation points, that can be modified to facilitate 
deeper and more productive usability testing. If the team failed to provide 
explicit variation capabilities for a suitable number of the implemented 
software's features, it would not fulfill the overall goal of the project. This 
could occur if sufficient planning and regular incremental development of the 
software's variation capabilities was not completed, as hapharzardly 
introducing feature variation to an already completed video game would not be 
able to fulfill the goal of the project.
\\

\noindent The aforementioned risks could compound to result in the required 
goals and scope of the RealmForge software not being fulfilled, or the software 
being delivered in an incomplete or erroneous state.
\\

\noindent To demonstrate that these risks can be overcome, the project's proof 
of concept demonstration will include the functionality outlined below. To 
ensure that the proof of concept build of the software is feasible to complete 
by the date of the demonstration, it is planned to be completely focused on 
providing the minimum functionality to address each risk. The proof of concept 
software build will include:
\begin{itemize}
	\item An initial version of the RealmForge software's main menu, 
	providing options to configure its feature variations and open the 
	primary combat game scene described below.
	\item An initial version of one of the software's enemy combat 
	encounter game scenes. This game scene will not include its planned 
	complete functionality, but it will include:
	\begin{itemize}
		\item A representation of the player character, that displays their 
		current health attribute
		\item A representation of an enemy (the player's opponent within 
		the encounter), that displays their current health attribute
		\item An initial version of the game scene's user interface, 
		providing buttons that correspond to the player's available actions 
		(these butttons will not 
		execute the complete planned game logic currently)
		\item An initial version of the game scene's combat functionality, 
		allowing simple actions selected by the user to damage the enemy 
		and 
		affect their health attribute. This could also be extended to 
		include the enemy responding with a simple action.
	\end{itemize}
	\item An initial version of the software's feature variation 
	capabilities. The following variation points will be demonstrated:
	\begin{itemize}
		\item A capability to modify the player's starting health (in the 
		final version of the software, users will be able to modify the 
		player's starting health to adjust the difficulty of their gameplay 
		sessions)
		\item A capability to toggle the combat encounter's initial user 
		interface described above between at least two configurations (a 
		different set or permutation of provided buttons)
	\end{itemize}
\end{itemize}

\noindent This planned scope will demonstrate that the RealmForge development 
team has 
the capabilities required to overcome the learning curve of their chosen game 
development tool (by completing a small version of the software, developed 
using the tool). It will also demonstrate that the project's need for explicit 
feature variation points can be fulfilled, and that the development team has 
the capabilities required to produce a functional software program within a 
feasible scope for the Capstone course.
\\

\noindent In the event of a failed proof of concept, the team has various 
methods of regrouping. First, the scope of the project could be pared down, 
focusing on more rigid plans and deadlines for core features. This would aid 
the direction of the project and ensure that it stays on course, avoiding scope 
creep through adding unnecessary features. Another method to refocus could be 
pivoting some key structure of the game to make it easier to fully flesh out. 
For example, changing the game engine/framework to something more familiar 
would allow the team to worry more about the quality of the work rather than 
learning while developing features. 

\section{Team Meeting Attendance}

\begin{table}[H]
\centering
\begin{tabular}{ll}
\toprule
\textbf{Student} & \textbf{Meetings}\\
\midrule
Total & 10\\
Andy Huynh  & 10\\
Andreas Sideris & 10\\
Humna Saeed & 10\\
Christopher Ludwig & 9\\
Owen Johnson & 10\\
\bottomrule
\end{tabular}
\end{table}

Christopher was excused from 1 team meeting, due to a valid absence.

\section{Supervisor/Stakeholder Meeting Attendance}

\noindent \textbf{Supervisor's Name: } Dr.~Charles Welch

\begin{table}[H]
\centering
\begin{tabular}{ll}
\toprule
\textbf{Student} & \textbf{Meetings}\\
\midrule
Total & 1\\
Andy Huynh  & 1\\
Andreas Sideris & 1\\
Humna Saeed & 0\\
Christopher Ludwig & 1\\
Owen Johnson & 1\\
\bottomrule
\end{tabular}
\end{table}

\noindent 1 preliminary meeting was held with the project's supervisor prior to 
the 
project's start, to finalize its details. All team members attended this 
meeting.

\section{Lecture Attendance}

\begin{table}[H]
\centering
\begin{tabular}{ll}
\toprule
\textbf{Student} & \textbf{Lectures}\\
\midrule
Total & 13\\
Andy Huynh  & 6\\
Andreas Sideris & 5\\
Humna Saeed & 0\\
Christopher Ludwig & 0\\
Owen Johnson & 6\\
\bottomrule
\end{tabular}
\end{table}

There were 13 lectures in total. 2 of them were TRON lectures, and 5 of them 
were from before issue tracking was mandatory/set up, which leaves 6 that 
attendance can be confirmed for. The team was not aware that lecture issues 
were required to be created prior to September 23rd, as this information was 
only provided in a line of the course navigation document that was missed.\\

\noindent \textbf{Note:} All team members (including Christopher and Humna) 
attended each 
lecture that was held before tracking began, except for solely Mechatronics 
focused lectures.

\section{TA Document Discussion Attendance}

\noindent \textbf{TA's Name: }Rashad Bhuiyan

\begin{table}[H]
\centering
\begin{tabular}{ll}
\toprule
\textbf{Student} & \textbf{Meetings}\\
\midrule
Total & 4\\
Andy Huynh  & 4\\
Andreas Sideris & 4\\
Humna Saeed & 3\\
Christopher Ludwig & 4\\
Owen Johnson & 4\\
\bottomrule
\end{tabular}
\end{table}

Humna was excused from 1 meeting for a valid reason.

\section{Commits, Lines Added and Lines Deleted, Pull Requests}

\begin{table}[H]
\centering
\begin{tabular}{lllllll}
\toprule
\textbf{Student} & \textbf{Commits} & \textbf{Percent} & \textbf{Lines Added} &
\textbf{Lines Deleted} & \textbf{Pull Requests}\\
\midrule
Total & 479 & 100\% & 26602 & 11486 & 110\\
Andy Huynh & 178 & 37\% & 17126 & 6685 & 19\\
Andreas Sideris & 118 & 25\% & 5092 & 3247 & 36\\
Humna Saeed & 52 & 11\% & 1727 & 722 & 16\\
Christopher Ludwig & 67 & 14\% & 896 & 345 & 16\\
Owen Johnson & 64 & 13\% & 1761 & 486 & 23\\
\bottomrule
\end{tabular}
\end{table}
Andy was responsible for forking the template repository and providing GitHub 
properties such as extra GitHub Actions and Issue Templates. This is why his line 
count and commit history is much higher than the rest. Roughly 12000 lines added and +50 commits. 
\section{Issue Tracker}

\begin{table}[H]
\centering
\begin{tabular}{lll}
\toprule
\textbf{Student} & \textbf{Authored (O+C)} & \textbf{Assigned (C only)}\\
\midrule
Andy Huynh & 31 & 87 \\
Andreas Sideris & 65 & 86 \\
Humna Saeed & 21 & 63 \\
Christopher Ludwig & 25 & 78 \\
Owen Johnson & 43 & 92 \\
\bottomrule
\end{tabular}
\end{table}

\noindent Some issues were created regarding multiple features at once, whereas 
others were more granular. 

\section{CICD via GitHub Actions}

% \wss{Say how CICD technology is used and will be used in your project.}

% \wss{Provide links to your CICD yaml files}
\noindent The current CI/CD checks for all major LaTeX files for correct file structure listed 
on the script/python documents (\textbf{select\_docs.py, buid.py}).
Once the CI/CD finishes checking for correctness, it builds the LaTeX files and deploys
them onto the GitHub page \href{https://aeolus-software.github.io/RealmForge/}{website} as PDFs downloadable.
See the \href{https://github.com/Aeolus-Software/RealmForge/blob/main/.github/workflows/latex-pages.yml}{latex-page.yml} on the repository.
\\
\noindent The CI/CD will provide a syntax linting and unit testing for the code in the 'src' file. These will 
help with catching early major errors and faults.
See the \href{https://github.com/Aeolus-Software/RealmForge/blob/main/.github/workflows/godot.yml}{godot.yml} on the repository.
\\
\noindent The CI/CD will provide a online Markdown version for the list of requirements, tests, and use cases such that
traceability and linking is available without access between the PDF documents. 
See the \href{https://github.com/Aeolus-Software/RealmForge/blob/main/.github/workflows/requirements.yml}{requirements.yml} on the repository.



\section{Custom Documents (CDs)}

The team will complete the following two CDs:

\begin{itemize}
	\item \textbf{CD-Fall:} Feature-Oriented Domain Analysis Report - initial 
	revision completed as of November 18. This document was suggested by 
	Dr.~Smith, as it is particularly relevant to the project due to the 
	software's 
	requirements for explicit feature variation points.
	\item \textbf{CD-Winter:} Usability Report - this document was selected 
	since 
	usability testing is a central focus of the project. It is still planned as 
	the project's CD-Winter.
\end{itemize}

\section{Team Charter Trigger Items}

No specific trigger items are outlined in the team charter as of this point. 
All team members have generally attended all (or almost all) required meetings, 
and completed their work to an acceptable standard.\\

\noindent To provide further accountability for team members, a new trigger 
item will be 
added to the team charter, to penalize team members who miss team meetings or 
TA meetings without a valid excuse. Team members violating this trigger item 
will have to purchase and provide refreshments/snacks for the next team meeting.

\section{Additional Productivity Metrics}

No other additional productivity metrics are collected currently.\\

% \wss{If your team has additional metrics of productivity, please feel free to
% add them to this report.  If not, please explicitly state that there are no additional
% metrics.}

% \wss{Additional metrics can include things like code reviews done, pull requests
% created, count of joining meetings late, count of number of times contributions
% had to be corrected, number of internal deadlines missed, test cases written, etc.}

% \wss{We are looking for data on these metrics, not just a list of additional
% metrics the team is planning on tracking.  However, if all you have is a plan,
% please share it here.}

\end{document}